% Clase 13 --- Por qué el promedio deja sólo la deriva (§1.2).
% El docente insistió en que el electrón "arranca con alguna velocidad, con una
% dirección aleatoria". La figura sigue la componente de la velocidad a lo largo
% del campo: en cada choque salta a un valor al azar (que puede ser negativo) y
% entre choque y choque crece con pendiente constante a = eE/m. Al promediar,
% los saltos al azar se cancelan y sobrevive sólo lo que aportó el campo.
\begin{center}
\begin{tikzpicture}
  \begin{axis}[
      fisfig,
      width=9.0cm, height=5.8cm,
      grid=none,
      axis lines=middle,
      xmin=-0.15, xmax=8.1, ymin=-1.35, ymax=1.95,
      xlabel={$t$}, ylabel={$v$ a lo largo de $\vec E$},
      xlabel style={anchor=north west}, ylabel style={anchor=south east},
      xtick=\empty, ytick={0}, yticklabels={$0$},
      axis line style={figgray},
    ]
    % tramos entre choques: pendiente constante a = eE/m
    \addplot[figblue, line width=1pt] coordinates {(0,0.6) (1.0,1.6)};
    \addplot[figblue, line width=1pt] coordinates {(1.0,-0.5) (2.2,0.7)};
    \addplot[figblue, line width=1pt] coordinates {(2.2,0.3) (3.1,1.2)};
    \addplot[figblue, line width=1pt] coordinates {(3.1,-0.7) (4.4,0.6)};
    \addplot[figblue, line width=1pt] coordinates {(4.4,0.2) (5.2,1.0)};
    \addplot[figblue, line width=1pt] coordinates {(5.2,-0.3) (6.5,1.0)};
    \addplot[figblue, line width=1pt] coordinates {(6.5,0.5) (7.5,1.5)};

    % los choques
    % OJO: sin \foreach --- dentro de un `axis` de pgfplots revienta con
    % este preámbulo (babel español); va desenrollado a mano.
    \node[figcarga, fill=figgray, minimum size=3.4pt] at (axis cs:1.0,0) {};
    \draw[figgray, dashed, line width=0.4pt] (axis cs:1.0,-0.95) -- (axis cs:1.0,1.7);
    \node[figcarga, fill=figgray, minimum size=3.4pt] at (axis cs:2.2,0) {};
    \draw[figgray, dashed, line width=0.4pt] (axis cs:2.2,-0.95) -- (axis cs:2.2,1.7);
    \node[figcarga, fill=figgray, minimum size=3.4pt] at (axis cs:3.1,0) {};
    \draw[figgray, dashed, line width=0.4pt] (axis cs:3.1,-0.95) -- (axis cs:3.1,1.7);
    \node[figcarga, fill=figgray, minimum size=3.4pt] at (axis cs:4.4,0) {};
    \draw[figgray, dashed, line width=0.4pt] (axis cs:4.4,-0.95) -- (axis cs:4.4,1.7);
    \node[figcarga, fill=figgray, minimum size=3.4pt] at (axis cs:5.2,0) {};
    \draw[figgray, dashed, line width=0.4pt] (axis cs:5.2,-0.95) -- (axis cs:5.2,1.7);
    \node[figcarga, fill=figgray, minimum size=3.4pt] at (axis cs:6.5,0) {};
    \draw[figgray, dashed, line width=0.4pt] (axis cs:6.5,-0.95) -- (axis cs:6.5,1.7);

    % el promedio
    \addplot[figred, line width=1.1pt, dashed, domain=0:7.5] {0.47};
    \node[figlbl, figred, anchor=west] at (axis cs:7.6,0.47) {$v_d$};

    \node[figlblaux, anchor=south] at (axis cs:5.85,-1.33) {choques};
  \end{axis}

  \node[figlbl, anchor=west, align=left] at (9.5,3.6)
    {\footnotesize arranca al azar:\\[-0.2em] \footnotesize promedia a $0$};
  \node[figlbl, figblue, anchor=west, align=left] at (9.5,2.2)
    {\footnotesize pendiente $a = \dfrac{eE}{m}$};
  \node[figlbl, figred, anchor=west, align=left] at (9.5,0.9)
    {\footnotesize lo único que\\[-0.2em] \footnotesize sobrevive};

\end{tikzpicture}\\[0.5em]
{\small\itshape Cada tramo es un movimiento uniformemente acelerado y cada punto
gris un choque, tras el cual el electrón arranca en una dirección cualquiera
---por eso los saltos van a valores positivos y negativos indistintamente---. La
deducción consiste en promediar $\mathbf v(t) = \mathbf v_0 - (e/m)\mathbf E\,t$
sobre muchísimos electrones: el término $\mathbf v_0$ se cancela porque las
direcciones de arranque no tienen nada que ver con el campo, y del término
restante queda $v_d = eE\tau/m$, con $\tau$ el tiempo libre medio. Notar el
orden de magnitud que la figura no puede respetar: los saltos verticales son del
orden de $10^{6}$ m/s y la línea roja está a algunos centímetros por hora, o sea
prácticamente pegada al eje. El paso decisivo del argumento es que
\textbf{$\tau$ casi no depende de $E$}: el movimiento está dominado por la
agitación térmica y el campo apenas lo sesga, así que la frecuencia de choques
no cambia ---y por eso $v_d$ sale proporcional a $E$, que es la ley de Ohm---.}
\end{center}
